\scanpage{13}
montrent que $f$ l'est aussi...], alors \unsure{le|la|là} \deleted{\unsure{d'}}on trouve
\[
 f'=\Delta f^{-1},\quad f''=\Delta' f'^{-1}=(\Delta'\Delta^{-1})f \tag*{(45)}
\]
ce qui s'écrit
\[
 \Delta f''=\Delta' f \tag*{(45 bis)}
\]
(d'après le ``principe de conservation des identités algébriques'', cette \deleted{\unsure{$f$}} identité reste vraie sans la condition \deleted{\unsure{sur}} $\Delta$ une unité).

De plus, dans le cas de validité de (45) ($\Delta$ unité) on en tire
\[
 \det(f')=\Delta^3\det(f)^{-1}=\Delta^3(\Delta\omega^{\unsure{\otimes}{-2}})^{-1}
\]
i.e.
\[
 \det f'=\Delta^2\omega^{\otimes2} \tag*{(\unsure{47|46})}
\]
\[
 \Delta'=\Delta^2\qquad\text{i.e.} \tag*{\deleted{(\unsure{48|47})}}
\]
\deleted{\unsure{qui reste}} de sorte que (45 bis) devient
\[
 f''=\Delta f \tag*{(48)}
\]
D'après le $=$ principe de conservation des identités Algébriques, cette identité reste valable pour \unsure{toute|tte|tout} $f$ (non seulement si $f$ est un isomorphisme). Ceci \deleted{\illegible{1}} \deleted{montre} .... que si $\Delta=1$ i.e. $f$ un iso et $\omega=\det f$, alors la relation entre $f$ et $f'$ est symétrique. Mais \unsure{nous|non},

% archive sheet number: 12
